% Results and conclusion text removed from main.tex on 2026-10-19 (AWGN numbers from the
% first stage-2 recipe; small budget, ablation and priority-embedding comparison).

% ============================================================
% IV. RESULTS
% ============================================================
\section{Results}
\label{sec:results}

\subsection{When Extra Tokens Help}
\label{ssec:saturation}

Without an allocator, sending four tokens per task instead of one adds
13 to 23 points of accuracy per task at $-10$\,dB on AWGN and 11 to 17
points on Rayleigh. The gain disappears from 5\,dB upward on AWGN,
while deep fades keep extra tokens useful up to 5\,dB on Rayleigh. With
the large budget, coming within 0.5 points of the best accuracy needs
29 tokens per task at $-15$\,dB but one from 5\,dB upward on AWGN, and
31 tokens at $-15$\,dB, 5 at 5\,dB and one only at 20\,dB on Rayleigh.
A fixed number of tokens therefore either wastes symbols on clean
channels or starves noisy ones.

\subsection{Small Budget}
\label{ssec:small}

With $K{=}8$ and $\lambda{=}0$, every method sends about 32 symbols and
the methods differ only in how they split the eight tokens. The learned
allocator is the best practical method, with $+0.29$ points over the
equal split on AWGN (95\% CI $[0.18, 0.41]$) and $+0.16$ points
($[0.03, 0.30]$) on Rayleigh, where the priority embedding does equally
well. The random split is worse than the equal split on both channels,
so how the tokens are split matters. The hindsight oracle, which knows
the labels and the noise, is 4.0 (AWGN) and 5.2 (Rayleigh) points above
the equal split.

With $K{=}12$ the cap no longer binds, and five prices
$\lambda \in \{0.001, 0.003, 0.01, 0.03, 0.1\}$ trace accuracy against
symbols. For the same accuracy, the allocator needs 21 to 37 percent
fewer symbols than fixed splits on AWGN and 9 to 19 percent fewer on
Rayleigh. The savings come from following the SNR
(Fig.~\ref{fig:persnr}): at $\lambda{=}0.01$ the allocator sends 46.5
symbols at $-15$\,dB and 12.0 at 20\,dB on AWGN (48.0 and 15.3 on
Rayleigh), against a constant 32 for the 3/3/2 split. At low SNR it
spends more and gains 4 to 5 points; from 10\,dB upward it sends 12 to
19 symbols and stays within 0.4 points of the fixed split.

\begin{figure}[!t]
  \centering
  \includegraphics[width=\columnwidth]{per_snr.pdf}
  \caption{Symbols sent (left) and mean accuracy (right) against SNR,
           $J{=}4$, $K{=}12$, $\lambda{=}0.01$, mean of 3 seeds. The
           fixed 3/3/2 split always sends 32 symbols.}
  \label{fig:persnr}
\end{figure}

\subsection{Large Budget}
\label{ssec:large}

The validation rule selects $\lambda{=}0.005$ on both channels. On the
test set, the allocator at $K{=}16$ sends 46 percent fewer symbols than
the equal split on AWGN (34.4 against 64) for $-0.11$ points
($[-0.21, -0.01]$), which passes the non-inferiority test. On Rayleigh
it sends 26 percent fewer (47.1 against 64) for $-0.24$ points
($[-0.35, -0.12]$); this price sat exactly at the margin on validation,
and the test misses non-inferiority by 0.1 points. One price lower,
the Rayleigh allocator saves 15 percent for $-0.04$ points. Sending all
384 symbols is more accurate, by 2.88 (AWGN) and 4.89 (Rayleigh)
points, almost entirely at $-15$ and $-10$\,dB, where the cap of 16
limits both the equal split and the allocator. From 10\,dB upward all
three systems are within 0.3 points, and the allocator sends 12 to 13
symbols on AWGN and 17 to 29 on Rayleigh instead of 64. On both
channels the allocator lies above the no-allocator curve at every
number of symbols (Fig.~\ref{fig:tradeoff}), and the saving grows with
the cap: 7, 46, 65 and 81 percent at caps 8, 16, 32 and 64 on AWGN,
and 1, 26, 49 and 72 percent on Rayleigh, at a cost of 0.6 to 2.0
points at the two largest caps.

\begin{figure}[!t]
  \centering
  \includegraphics[width=\columnwidth]{tradeoff.pdf}
  \caption{Mean accuracy against symbols per sentence, $J{=}32$, test
           set. Grey: no allocator, every task sends $k$ tokens. Blue:
           the allocator at $K{=}16$ for each price $\lambda$; the
           circle marks the price selected on validation.}
  \label{fig:tradeoff}
\end{figure}

\subsection{Allocator Inputs}
\label{ssec:ablation}

We train ten allocators with different inputs at $K{=}16$ and the
selected $\lambda$, two of them with an extra confidence input (the
normalized entropy of each head on noise-free tokens), and evaluate
them at equal, ESG-heavy and Sentiment-heavy priorities on the
objective \eqref{eq:problem} and the weighted accuracy. On AWGN every
input matters: without the SNR the allocator sends about 51 symbols at
every SNR, and removing the priority or the sentence summary worsens
the objective by 0.020 ($[0.017, 0.023]$) and 0.027
($[0.024, 0.030]$). On Rayleigh the SNR and the priority matter
(removing the priority costs 0.039, $[0.035, 0.043]$), but the sentence
summary does not help: without it the objective is 0.007 lower
($[0.005, 0.009]$). The confidence input does not help on either
channel.

\subsection{Priority Steering}
\label{ssec:priority}

We raise the ESG priority from 0.1 to 0.9, with the other two tasks
sharing the rest equally, and compare the gain form \eqref{eq:score}
with the priority embedding, with the small budget at 0\,dB ($K{=}8$)
and the large budget at $-10$ to 0\,dB ($K{=}16$). In all four
settings, the gain form moves more tokens to ESG and raises ESG
accuracy more, by a factor of 1.7 to 22: with the large budget, by
5.76 against 2.55 points on AWGN and 8.27 against 4.99 points on
Rayleigh. With the small budget, the embedding gives fewer ESG tokens
to 0.03 percent (AWGN) and 2.2 percent (Rayleigh) of the sentences
when the ESG priority rises, while the gain form never does, as
Proposition~\ref{prop:monotone} guarantees. Compared with an allocator
trained only at $\mathbf{w} = (0.1, 0.1, 0.8)$, the single model
trained with random priorities is equivalent on AWGN ($-0.08$ points,
$[-0.31, 0.17]$) and 0.58 points worse on Rayleigh ($[0.11, 1.06]$).

% ============================================================
% V. CONCLUSION
% ============================================================
\section{Conclusion}
\label{sec:conclusion}

We moved the split between tasks in multi-task semantic communication
to the transmitter. Each task owns ordered task tokens, and a small
allocator decides for every sentence how many of them each task sends,
from the sentence, the average SNR and the priority, under a cap and a
price per token, with a priority input that provably never works
against the chosen task. The allocator follows the channel, splits a
fixed budget better than fixed rules, reaches the accuracy of fixed
splits with clearly fewer symbols, and steers accuracy toward the
prioritized task more strongly than a standard priority embedding. The
FLS and ESG labels come from classifiers, the training SNR range does
not cover $-15$\,dB, and the transmitter knows the average SNR without
error. Future work will consider imperfect channel reports, a dataset
with human labels for every task, and frequency-selective channels,
where the transmitter could also choose which channel uses carry each
task.

